\chapter{Hiện Thực Hóa RTL Đơn Nhân}
\label{chap:implementation_core}

Chương này trình bày chi tiết việc hiện thực hóa các thành phần của lõi vi xử lý RISC-V 5 tầng pipeline (\texttt{risc\_core}) ở mức RTL (Register Transfer Level). Mỗi module quan trọng sẽ được phân tích về mục đích, giao tiếp, thuật toán/FSM, kèm theo đoạn mã tiêu biểu, sơ đồ khối và biểu đồ thời gian.

\section{Tổ chức mã nguồn lõi xử lý}
Toàn bộ mã nguồn RTL của lõi được tổ chức trong thư mục \texttt{rtl/core/}, tuân thủ nguyên tắc thiết kế module hóa.
\begin{itemize}
    \item \texttt{risc\_core.v}: Module đỉnh của lõi, ghép nối datapath, hazard controller và giao tiếp bộ nhớ.
    \item \texttt{alu.v}: Đơn vị tính toán số học logic (ALU).
    \item \texttt{control\_unit.v}: Bộ điều khiển trung tâm, giải mã opcode.
    \item \texttt{reg\_file.v}: Tập 32 thanh ghi đa năng.
    \item \texttt{pc\_logic.v}: Logic tính toán cập nhật Program Counter.
    \item \texttt{pipeline\_register/}: Chứa 4 thanh ghi ranh giới (\texttt{if\_id.v}, \texttt{id\_ex.v}, \texttt{ex\_mem.v}, \texttt{mem\_wb.v}).
    \item \texttt{instr\_mem.v}: Bộ nhớ lệnh (Instruction Memory).
    \item \texttt{data\_mem.v}: Bộ nhớ dữ liệu (Data Memory).
\end{itemize}

\section{Bộ nhớ và cơ chế nạp chương trình}
Trong hệ thống thực tế, bộ nhớ lệnh (Instruction Memory) và bộ nhớ dữ liệu (Data Memory) thường là các vi mạch SRAM hoặc Flash. Trong môi trường mô phỏng RTL, chúng được mô hình hóa bằng các mảng thanh ghi (array of registers) trong Verilog. 

Để nạp chương trình vào bộ nhớ lệnh mà không cần phải mã hóa cứng từng dòng dữ liệu, dự án sử dụng hàm hệ thống \texttt{\$readmemh} (Read Memory Hex) của Verilog. Hàm này cho phép trình biên dịch đọc nội dung từ một tệp văn bản chứa các mã lệnh dưới dạng hệ thập lục phân (hexadecimal) và khởi tạo trực tiếp vào mảng bộ nhớ.
\begin{lstlisting}[caption={Cơ chế nạp file .hex vào bộ nhớ},label={lst:readmemh}]
reg [31:0] memory [0:255]; // Mang bo nho 256 tu 32-bit

initial begin
    // Nap noi dung tu file "program.hex" vao mang memory
    $readmemh("program.hex", memory);
end
\end{lstlisting}
Khi mô phỏng bắt đầu, khối \texttt{initial} sẽ được thực thi đầu tiên. Tệp \texttt{program.hex} được tạo ra từ việc biên dịch mã nguồn Assembly (hoặc C/C++) bằng công cụ RISC-V GNU Toolchain, sau đó trích xuất dưới dạng văn bản plain-text. Máy tính mô phỏng đọc từng dòng văn bản hex, chuyển đổi thành giá trị nhị phân và gán vào các phần tử tương ứng của mảng \texttt{memory}, từ đó lõi vi xử lý có thể fetch các lệnh này qua cổng \texttt{instr\_from\_mem}.

\section{Tập thanh ghi đa năng (Register File)}
\textbf{Mục đích:} Lưu trữ 32 thanh ghi 32-bit theo chuẩn RV32I.
\textbf{Giao tiếp:} Hai cổng đọc tổ hợp (\texttt{rs1\_addr}, \texttt{rs2\_addr}) và một cổng ghi đồng bộ (\texttt{rd\_addr}, \texttt{wdata}, \texttt{we}).
\textbf{Đặc điểm hiện thực:} Thanh ghi \texttt{x0} được nối cứng bằng 0. Việc đọc được thực hiện bất đồng bộ để đảm bảo dữ liệu sẵn sàng ngay trong chu kỳ Decode.

\begin{lstlisting}[caption={Hiện thực Register File},label={lst:reg_file}]
always @(posedge clk) begin
    if (we && rd_addr != 5'd0) begin
        registers[rd_addr] <= wdata;
    end
end
assign rs1_data = (rs1_addr == 5'd0) ? 32'd0 : registers[rs1_addr];
assign rs2_data = (rs2_addr == 5'd0) ? 32'd0 : registers[rs2_addr];
\end{lstlisting}

\section{Đơn vị số học logic (ALU)}
\textbf{Mục đích:} Thực thi các phép toán số học (cộng, trừ), logic (AND, OR, XOR) và dịch bit.
\textbf{Giao tiếp:} Nhận hai toán hạng 32-bit \texttt{a}, \texttt{b} và mã điều khiển \texttt{alu\_control[3:0]}. Trả về \texttt{result} 32-bit và cờ \texttt{zero}.

\section{Bộ điều khiển trung tâm (Control Unit)}
\textbf{Mục đích:} Giải mã chỉ thị 32-bit, phát ra cờ điều khiển cho ALU, nhánh, và giao tiếp bộ nhớ.
\textbf{Thuật toán/Logic:} Sử dụng câu lệnh \texttt{case} trên trường \texttt{opcode[6:0]} và \texttt{funct3}, \texttt{funct7} để thiết lập trạng thái. Nó cũng sinh ra giá trị mở rộng dấu \texttt{imm} phù hợp với định dạng lệnh.

\section{Thiết kế Datapath và Pipeline 5 tầng}
Đường truyền dữ liệu được chia làm 5 tầng: Fetch, Decode, Execute, Memory, Writeback. Các thanh ghi pipeline đóng vai trò chốt dữ liệu và tín hiệu điều khiển giữa các tầng.

\subsection{Tầng Fetch và IF/ID}
Bao gồm bộ đếm PC và mạch \texttt{pc\_logic}. PC được cập nhật dựa trên tín hiệu rẽ nhánh từ tầng EX hoặc stall từ hazard controller. Thanh ghi IF/ID chốt mã lệnh từ \texttt{instr\_mem}.

\subsection{Tầng Decode và ID/EX}
Thực hiện giải mã qua \texttt{control\_unit} và đọc \texttt{reg\_file}. Khi gặp stall, tín hiệu \texttt{id\_ex\_valid} bị ép về 0 (tạo NOP).

\subsection{Tầng Execute và Forwarding}
Chứa ALU và mạch so sánh rẽ nhánh. Mạch Forwarding ưu tiên lấy dữ liệu mới nhất từ MEM hoặc WB để thay thế dữ liệu từ ID/EX nhằm tránh data hazard.

\begin{table}[H]
\centering
\begin{tabular}{|l|c|c|c|c|c|c|}
\hline
\textbf{Lệnh} & \textbf{C1} & \textbf{C2} & \textbf{C3} & \textbf{C4} & \textbf{C5} & \textbf{C6} \\
\hline
1: \texttt{ADD x3, x1, x2} & IF & ID & EX & MEM & \textbf{WB} & \\
\hline
2: \texttt{SUB x4, x3, x5} & & IF & ID & \textbf{EX*} & MEM & WB \\
\hline
\end{tabular}
\caption{Biểu đồ thời gian minh họa cơ chế Forwarding (Từ EX/MEM về EX của lệnh sau)}
\label{fig:forwarding_wave}
\end{table}

\subsection{Tầng Memory và MEM/WB}
Giao tiếp với bus hệ thống. Nếu \texttt{mem\_ready=0}, lõi kéo tín hiệu \texttt{mem\_wait=1} để đóng băng pipeline (stall).

\begin{table}[H]
\centering
\begin{tabular}{|l|c|c|c|c|c|c|}
\hline
\textbf{Lệnh} & \textbf{C1} & \textbf{C2} & \textbf{C3} & \textbf{C4} & \textbf{C5} & \textbf{C6} \\
\hline
1: \texttt{LW x3, 0(x1)} & IF & ID & EX & \textbf{MEM} (Wait) & \textbf{MEM} (Ready) & WB \\
\hline
2: \texttt{ADD x4, x3, x5} & & IF & ID & \textbf{EX} (Stall) & \textbf{EX} (Stall) & EX \\
\hline
\end{tabular}
\caption{Biểu đồ thời gian minh họa Memory Stall khi truy cập RAM bị trễ}
\label{fig:mem_stall_wave}
\end{table}

\section{Xử lý xung đột (Hazard Controller)}
\textbf{Mục đích:} Phát hiện và xử lý Data Hazard (bằng Forwarding và Load-Use Stall) và Control Hazard (bằng Branch Flush).

\subsection{Load-Use Stall}
Xảy ra khi lệnh EX là lệnh Load và có thanh ghi đích trùng với nguồn của lệnh ID.
\begin{table}[H]
\centering
\begin{tabular}{|l|c|c|c|c|c|c|c|}
\hline
\textbf{Lệnh} & \textbf{C1} & \textbf{C2} & \textbf{C3} & \textbf{C4} & \textbf{C5} & \textbf{C6} & \textbf{C7} \\
\hline
1: \texttt{LW x3, 0(x1)} & IF & ID & EX & MEM & \textbf{WB} & & \\
\hline
2: \texttt{ADD x4, x3, x5} & & IF & ID & \textbf{Stall} & \textbf{EX*} & MEM & WB \\
\hline
\end{tabular}
\caption{Biểu đồ thời gian minh họa Load-Use Stall (Chèn 1 chu kỳ NOP để đợi Load)}
\label{fig:load_use_wave}
\end{table}

\subsection{Branch Flush}
Khi nhánh được xác định là rẽ (Taken) ở tầng EX, hai lệnh đi sau ở IF và ID bị hủy (Flush).
\begin{table}[H]
\centering
\begin{tabular}{|l|c|c|c|c|c|c|}
\hline
\textbf{Lệnh} & \textbf{C1} & \textbf{C2} & \textbf{C3} & \textbf{C4} & \textbf{C5} & \textbf{C6} \\
\hline
1: \texttt{BEQ x1, x2, Target} & IF & ID & \textbf{EX} (Taken) & MEM & WB & \\
\hline
2: \texttt{Lệnh kế tiếp} & & IF & ID & \textbf{Flush} & & \\
\hline
3: \texttt{Lệnh kế tiếp 2} & & & IF & \textbf{Flush} & & \\
\hline
4: \texttt{Lệnh tại Target} & & & & \textbf{IF} & ID & EX \\
\hline
\end{tabular}
\caption{Biểu đồ thời gian minh họa Branch Flush (Hủy bỏ 2 lệnh sai luồng)}
\label{fig:branch_flush_wave}
\end{table}
